\section{Future work}
In its current state, the system is able to generate environments containing vegetation and rivers, which can be manipulated using a set of parameters. However, there are still some improvements which can be made to the system to improve the features and capabilities of the system. A few of these improvements are described in the following section.

\subsection{Improving performance}
In the initial version of the system, there has been no focus on improving the performance of the system. One improvement that has been made is caching the intersection between fat curve triangles and chunks, but the performance of the system can be further improved by taking advantage of Unity's systems. One example would be to use parallelisation and SIMD by using Unity's job system in generating the terrain or operations with fat curves. A paper by Wei \cite{wei_parallel_2008} proposed a way to parallelise the Poison Disk Distribution, which can be applied to the object placement.

Although the current system is still able to generate landscapes in a reasonable time, shorter execution times can allow the user to generate more terrains in the same timespan, allowing for more choice in the resulting environment. \todo{More reasons?}

\subsection{Multiple environments}
While the system is able to generate environments with multiple types of vegetation, is not able to generate terrains with multiple types of environments. For example, the system cannot generate the edge of forests where a section of the terrain contains trees and the other section only contains grass or bushes. Interesting environments often contain different types of landscape, so, therefore, the addition of multiple types of environment should be added. Some research has already gone into the generation of multiple ecotope. Hammes \cite{hammes_modeling_2001} uses terrain features such as the height and slope of the terrain to determine which ecotope the section of terrain belongs to. Based on the ecotope, different plants are selected for placement. \cite{do_nascimento_gpu-based_2018}  and \cite{van_muijden_gpu-based_2017} take a similar approach, but both use water as well as the terrain, such that rivers also influence the selected ecotype.

\subsection{Generating roads}
In digital nature research, the presence of artificial objects is a desired component of the system \cite{van_houwelingen-snippe_designing_2023}. Because of this, one of the next features that should be added is the presence of roads or road networks. Roads can be planned on the terrain, and the addition will change the architecture. The planning and generation of roads has been researched before, most times in the procedural generation of cityscapes \cite{kelly_citygen_nodate} or in the procedural generation system by Smelik \cite{smelik_declarative_2011}. In addition to the placement of the roads itself, certain objects can be placed adjacent to roads, such as roads, sign posts or bins.

\subsection{Generating from point of view}
Although the aim of the system is to generate an environment for a single point of view, the generation techniques used are the same that are used to create large-scale environments. New techniques that involve the position of the camera more directly to generate each component can generate the environment more effectively. As mentioned in the discussion, rivers are not likely to be generated visible to the camera. When the position of the camera is taken into account, rivers can be generated such that they will appear in a certain position of the view. Additionally, the generation of objects can be changed by including the position of the camera to distinguish between foreground and background objects, which can result in interesting environments.
